\documentclass[11pt,a4paper]{article}

\usepackage[utf8]{inputenc}
\usepackage[T1]{fontenc}
\usepackage{lmodern}
\usepackage{amsmath,amssymb,amsfonts}
\usepackage{booktabs}
\usepackage{hyperref}
\usepackage[margin=2.5cm]{geometry}
\usepackage{graphicx}
\usepackage{xcolor}
\usepackage{authblk}

\hypersetup{
  colorlinks=true,
  linkcolor=blue!60!black,
  citecolor=blue!60!black,
  urlcolor=blue!60!black
}

\title{\textbf{Lattice Sigma Terms as an Anchor\\for the Dense Nuclear Matter Equation of State}}

\author{Oleg Kirichenko}
\affil{Independent Researcher; \texttt{urevich55@gmail.com}}

\date{\today}

\begin{document}

\maketitle

\begin{abstract}
A phenomenological Vacuum Mass Fraction (VMF) model is proposed, in which the hadron mass is separated into two components: (1)~the current mass, determined by the explicit breaking of chiral symmetry via sigma terms, and (2)~the residual nonperturbative component $M_\Omega$, generated by confinement, gluon field energy, and the QCD trace anomaly.

Using lattice QCD data for the pion-nucleon ($\sigma_{\pi N} \approx 44$~MeV) and strange ($\sigma_{sN} \approx 30$~MeV) sigma terms, the vacuum value of the nonperturbative component is fixed at $M_{\Omega,0} = 859 \pm 8$~MeV, constituting ${\sim}\,91\%$ of the nucleon mass.

A series of EOS prototypes for dense nuclear matter is constructed with a two-parameter ansatz for the in-medium modification $M_\Omega(n_B)$. It is shown that: (i)~minimal realizations with constant vector repulsion yield a superluminal speed of sound; (ii)~density-dependent saturation of the vector field resolves the causality problem ($c_s^2 \leq 1$); (iii)~the problem of an overly stiff equation of state is resolved by introducing a first-order phase transition to the conformal QGP limit, reducing the maximum neutron star mass to a realistic $M_{\max} \approx 2.3\,M_\odot$.

The model is strictly bounded by input data: the vacuum scale $M_{\Omega,0}$ is fixed by lattice calculations and is not a free fitting parameter.

\medskip
\noindent\textbf{Keywords:} hadron mass, sigma terms, equation of state, neutron stars, nonperturbative QCD
\end{abstract}

\tableofcontents

%======================================================================
\section{Introduction}
\label{sec:intro}

\subsection{The Nucleon Mass Problem}

The proton mass $M_N \approx 938.3$~MeV is only to a small extent determined by the current quark masses. According to the Ji decomposition~\cite{Ji1995}, QCD sum rules~\cite{SVZ1979}, and the analysis of the energy-momentum tensor trace, up to ${\sim}\,90\%$ of the mass is generated by nonperturbative mechanisms: gluon field energy and the trace anomaly.

At the operator level:
\begin{equation}
T^{\mu}{}_{\mu} = \frac{\beta(g)}{2g}\, G^a_{\alpha\beta} G^{a\alpha\beta} + \bigl(1+\gamma_m(g)\bigr) \sum_q m_q \bar{q} q\,.
\label{eq:trace}
\end{equation}

Even in the chiral limit ($m_q \to 0$), the gluon term remains---mass is generated nonperturbatively.

\subsection{Aim of the Work}

We propose using standard QCD sigma terms to define a dimensionless parameter $f_\Omega$, which characterizes the fraction of the hadron mass not reducible to current quark masses, and to investigate the behavior of this parameter in a dense baryon medium through its impact on the equation of state of nuclear matter.

%======================================================================
\section{Mass Decomposition via Sigma Terms}
\label{sec:decomposition}

\subsection{Definitions}

Via the Feynman--Hellmann theorem, the standard sigma terms are:
\begin{equation}
\sigma_{qH} = m_q \frac{\partial M_H}{\partial m_q}\,.
\end{equation}

The residual nonperturbative mass of hadron $H$:
\begin{equation}
\boxed{\;M_\Omega^{(H)}[\mathcal{F}] \equiv M_H - \sum_{q \in \mathcal{F}} \sigma_{qH}\;}
\end{equation}

The dimensionless fraction:
\begin{equation}
\boxed{\;f_\Omega^{(H)}[\mathcal{F}] = 1 - \frac{1}{M_H} \sum_{q \in \mathcal{F}} \sigma_{qH}\;}
\end{equation}

\subsection{Evaluation for the Nucleon}

Using lattice data~\cite{Agadjanov2023,Gupta2021}:

\begin{table}[h]
\centering
\begin{tabular}{lll}
\toprule
Quantity & Value & Source \\
\midrule
$\sigma_{\pi N}$ & $44 \pm 3$~MeV & Gupta et al.\ (2021)~\cite{Gupta2021} \\
$\sigma_{sN}$ & $30 \pm 7$~MeV & Agadjanov et al.\ (2023)~\cite{Agadjanov2023} \\
$\sigma_{\text{heavy}}$ & ${\sim}\,6$~MeV & PDG matching~\cite{PDG2024} \\
$\sum_q \sigma_{qN}$ & $80 \pm 15$~MeV & Sum \\
\bottomrule
\end{tabular}
\caption{Sigma-term input data for the nucleon.}
\label{tab:sigma}
\end{table}

We obtain:
\begin{equation}
f_\Omega^{(N)} \approx 1 - \frac{80 \pm 15}{939} \approx 0.91 \pm 0.02\,,
\qquad
M_{\Omega,0} = 859 \pm 8~\text{MeV}\,.
\end{equation}

Within this phenomenological decomposition, about $90\%$ of the nucleon mass is associated with nonperturbative QCD contributions not reducible to explicit current quark masses.

\subsection{Connection to the Ji Decomposition}

In the standard decomposition:
\begin{equation}
M_H = M_H^{\text{q,kin}} + M_H^{\text{g}} + M_H^{\text{q,mass}} + M_H^{\text{anom}}\,,
\end{equation}
our residual mass corresponds to:
\begin{equation}
M_\Omega^{(H)} \sim M_H^{\text{g}} + M_H^{\text{anom}} + M_H^{\text{conf}}\,.
\end{equation}

This is a phenomenological identification, not a theorem (cf.\ bag models~\cite{Chodos1974}). Its value lies in organization: a single variable $f_\Omega$ compresses the question of ``what fraction of mass does not reduce to current masses'' into a single number.

\subsection{Status of the Statement}

The quantity $f_\Omega^{(N)}$ is \emph{not} a new fundamental operator or conserved charge. It is an EFT parameter that organizes the hadron mass budget. Its value is stable upon fixing the flavor set $\mathcal{F}$ and is compatible with independent lattice decompositions.

%======================================================================
\section{Field Structure of the Model}
\label{sec:field}

The action of the VMF model has the standard form of GR + a canonical scalar field:
\begin{equation}
S = \int d^4x \sqrt{-g} \left[ \frac{c^4}{16\pi G}(R - 2\Lambda_{\text{eff}}) - \frac{1}{2}\partial_\mu \mathcal{W} \partial^\mu \mathcal{W} - V(\mathcal{W}) + \mathcal{L}_m \right].
\label{eq:action}
\end{equation}

The Einstein tensor is unmodified, and the causal cone is preserved. The construction is mathematically equivalent to standard scalar-tensor models with conformal coupling $\omega \to \infty$ (the Brans--Dicke limit~\cite{Will2014}), guaranteeing $\gamma_{\text{PPN}} = 1$ and automatic compatibility with Solar System constraints (Cassini: $|\gamma - 1| < 2.3 \times 10^{-5}$~\cite{Cassini2003}).

The vacuum state is defined as $\mathcal{W}_0$, where $V'(\mathcal{W}_0) = 0$, $V''(\mathcal{W}_0) > 0$. As $n_B \to 0$, all medium fields vanish, and the model is identical to classical GR.

%======================================================================
\section{Equation of State of Dense Matter}
\label{sec:eos}

\subsection{Scalar Field as an Order Parameter}

We postulate that the scalar field $\mathcal{W}$ acts as a macroscopic order parameter for the nonperturbative QCD condensate:
\begin{equation}
M_\Omega = g_\Omega \langle \mathcal{W} \rangle\,.
\end{equation}

In a dense baryon medium, the dynamics of the order parameter is determined by minimizing the effective potential:
\begin{equation}
V_{\text{eff}}(\mathcal{W}, n_B) = V(\mathcal{W}) + n_B\, M^*(\mathcal{W})\,.
\end{equation}

Since the exact form of $V(\mathcal{W})$ from first principles is unknown, we parameterize the \emph{solution} through a phenomenological ansatz:
\begin{equation}
M_\Omega(n_B) = M_{\Omega,0} \left(1 + \kappa_2 \frac{n_B}{n_0}\right)^{-\kappa_1/\kappa_2},
\label{eq:ansatz}
\end{equation}
where $\kappa_1, \kappa_2$ are phenomenological medium-response parameters, \emph{not derived} from first principles of QCD.

The current component depends on density through the chiral condensate:
\begin{equation}
M_{\text{current}}(n_B) = M_{\text{current},0} \cdot \frac{\langle \bar{q} q \rangle_{n_B}}{\langle \bar{q} q \rangle_0}\,.
\end{equation}

In the low-density approximation:
\begin{equation}
\frac{\langle \bar{q} q \rangle_{n_B}}{\langle \bar{q} q \rangle_0} \approx 1 - \frac{\sigma_{\pi N}\, n_B}{f_\pi^2\, m_\pi^2}\,,
\end{equation}
so the total effective nucleon mass takes the form:
\begin{equation}
M^*(n_B) \approx M_{\text{current},0} \left(1 - \frac{\sigma_{\pi N}\, n_B\, (\hbar c)^3}{f_\pi^2\, m_\pi^2}\right) + M_\Omega(n_B)\,.
\label{eq:mstar}
\end{equation}

\subsection{Energy Density}

\begin{equation}
\varepsilon(n_B) = \sum_{i=n,p,e,\mu} \varepsilon_{\text{kin}}^i + V(\langle \mathcal{W} \rangle) + \varepsilon_{\text{vec}}(n_B) + \varepsilon_{\rho}(n_B)\,,
\end{equation}
with density-dependent vector saturation:
\begin{equation}
\varepsilon_{\text{vec}}(n_B) = \frac{C_{\omega 0}}{2} \cdot g(n_B) \cdot n_B^2, \quad g(x) = \frac{1}{1 + \alpha_v(x-1)^{\nu_v}} \;\text{for}\; x > 1\,.
\end{equation}

\subsection{Calibration}

The model is calibrated by nuclear physics, \emph{not} pulsars:

\begin{table}[h]
\centering
\begin{tabular}{lll}
\toprule
Constraint & Value & What it fixes \\
\midrule
Binding energy & $E/A = -16$~MeV at $n_0$ & $C_{\omega 0}$ \\
Pressure & $P(n_0) = 0$ & Saturation \\
Symmetry energy & $S(n_0) = 32$~MeV & $C_\rho$ \\
Vacuum mass & $M_{\Omega,0} = 859 \pm 8$~MeV & \textbf{Fixed by lattice} \\
\bottomrule
\end{tabular}
\caption{Nuclear physics calibration constraints.}
\label{tab:calibration}
\end{table}

\subsection{The Double Counting Problem}

When merging the nonperturbative QCD condensate with macroscopic EOS, a risk of double counting arises. This is resolved by the fact that $M_{\Omega,0}$ acts exclusively as a \textbf{vacuum anchor}. In vacuum ($n_B = 0$), macroscopic medium fields are absent, and $M_{\Omega,0}$ strictly corresponds to the lattice value. The interaction constants of the macroscopic sector are \textbf{renormalized} around this rigid vacuum anchor during calibration (Sec.~\ref{sec:eos}). The total energy density is strictly tied to nuclear observables, eliminating physical energy duplication.

\subsection{Sensitivity Analysis}

The total uncertainty from first principles:
\begin{equation}
\Delta M_{\Omega,0} = \sqrt{(\Delta \sigma_{\pi N})^2 + (\Delta \sigma_{sN})^2} \approx \sqrt{3^2 + 7^2} \approx 7.6~\text{MeV}\,,
\end{equation}
which is less than $1\%$ of the baseline $M_{\Omega,0} = 859$~MeV. The maximum neutron star mass $M_{\max}$ varies within $\pm 0.05\,M_\odot$ under extreme sigma-term shifts ($1\sigma$ bounds), proving the robustness of the EFT predictions.

\subsection{Model Constraints}

\begin{enumerate}
\item Lattice QCD establishes $M_{\Omega,0} \in [851, 867]$~MeV.
\item Causality: $c_s^2 \le 1$ at all densities.
\item Astrophysical compatibility: $M_{\max} \geq 2.01\,M_\odot$, $R_{1.4} \in [11.5, 13.0]$~km.
\end{enumerate}

%======================================================================
\section{Computational Results}
\label{sec:results}

\subsection{Best Fit Point}

With the saturating vector field sector:

\begin{table}[h]
\centering
\begin{tabular}{ll}
\toprule
Parameter & Value \\
\midrule
$\kappa_1$ & 0.25 \\
$\kappa_2$ & 0.80 \\
$\alpha_v$ & 4.0 \\
$\nu_v$ & 2.0 \\
$M_D(n_0)/M_N$ & 0.675 \\
$c_{s,\max}^2$ & 0.393 \\
\bottomrule
\end{tabular}
\caption{Best-fit parameter set.}
\label{tab:bestfit}
\end{table}

\subsection{Resolving the Overly Stiff EOS}

The minimal model yielded $M_{\max} \sim 4.3\text{--}4.7\,M_\odot$, far exceeding the observational constraint $M_{\max} \lesssim 2.3\,M_\odot$. Since the fixed QCD input does not allow arbitrary parameter fitting, the model itself pointed to the necessity of a \textbf{first-order phase transition} to conformal QGP ($P = \varepsilon/3$) at $n_B \approx 2.0\,n_0$. Solving the TOV equations yields $M_{\max} \approx 2.3\,M_\odot$ and $R_{1.4} \approx 12$~km.

%======================================================================
\section{Extrapolation to the Black Hole Interior: A Heuristic Analysis}
\label{sec:bh}

\emph{This section is exploratory and lies outside the strict EFT framework of Sections~\ref{sec:decomposition}--\ref{sec:results}. All results are explicitly marked as heuristic.}

\subsection{Concept}

One speculative hypothesis is the extrapolation of $M_\Omega(n_B) \to 0$ to its logical limit inside the event horizon. Outside the horizon ($n_B \to 0$), the model is identically GR (Schwarzschild/Kerr). Inside, at $n_B \gg n_0$, the mass melts and matter transitions to conformal QGP:
\begin{equation}
P = \frac{1}{3}\varepsilon, \qquad c_s^2 = \frac{1}{3} < 1 \quad (\text{causal})\,.
\end{equation}

\subsection{Phenomenological Stitching}

The VMF model alone predicts $P = \varepsilon/3$, which satisfies the strong energy condition and \emph{inevitably leads to a singularity}. The transition to a quantum core ($P_r = -\varepsilon_{\max}$) requires physics beyond this EFT. We present a phenomenological stitching with Markov's limiting density hypothesis. This strategy is widely applied in the literature~\cite{Bardeen1968,Dymnikova1992,Hayward2006,Simpson2019}. Our contribution is the fixation of $r_0$ via QCD parameters rather than free fitting.

The heuristic density profile $\varepsilon(r) = \varepsilon_{\max} r_0^6 / (r^3 + r_0^3)^2$ yields the Hayward metric:
\begin{equation}
ds^2 = -\!\left(1 - \frac{2GM r^2}{r^3 + r_0^3}\right) dt^2 + \left(1 - \frac{2GM r^2}{r^3 + r_0^3}\right)^{-1}\! dr^2 + r^2 d\Omega^2\,.
\label{eq:hayward}
\end{equation}

As $r \to 0$, a regular de~Sitter core forms. This illustrates \emph{mathematical compatibility} of the QCD scale with regular metrics, not a dynamical prediction.

\subsection{Falsifiability}

\begin{enumerate}
\item If NICA/FAIR do not detect in-medium mass modifications at $n_B \sim 3\text{--}5\,n_0$, the model is falsified.
\item Gravitational-wave echoes from a regular core in post-merger signals would support the model.
\item The exterior geometry is identically Schwarzschild/Kerr; any deviations falsify the extrapolation.
\end{enumerate}

%======================================================================
\section{Observational Constraints}
\label{sec:obs}

\textbf{Gravity:}
LIGO-Virgo-KAGRA (O4a): 42 BH mergers, no deviations from GR~\cite{LVK2024}. Tensor mode speed: $|c_T/c - 1| \lesssim 10^{-15}$~\cite{Creminelli2017}. Cassini: $|\gamma_{\text{PPN}} - 1| < 2.3 \times 10^{-5}$~\cite{Cassini2003}.

\textbf{Hadron physics:}
Lattice QCD (sigma terms, gluon fractions); near-threshold $J/\psi$ photoproduction (scalar gluon structure); neutron stars via NICER and GW ($\kappa_\Omega(n_B)$---vacuum mass response to density).

%======================================================================
\section{Connection with Existing Dense Matter Models}
\label{sec:connection}

\textbf{Similarities:} Density-dependent effective mass and vector saturation are standard RMF elements~\cite{Walecka1974,Serot1986}. In-medium condensate modification relates to Brown--Rho scaling~\cite{BrownRho1991}. The conformal crossover is characteristic of quarkyonic matter and NJL models; see reviews~\cite{Lattimer2016,Drischler2021}.

\textbf{Key novelty:} Unlike RMF models where the nucleon mass is a fitting parameter, here $M_{\Omega,0}$ is strictly fixed by lattice QCD via sigma terms. The single organizing parameter $f_\Omega$ links nuclear physics (symmetry energy) with hadron phenomenology (vector meson mass shifts).

%======================================================================
\section{Numerical Implementation and Phenomenology}
\label{sec:numerics}

\begin{table}[h]
\centering
\begin{tabular}{llll}
\toprule
Hadron & Total Mass (MeV) & $M_\Omega$ (MeV) & $f_\Omega$ (\%) \\
\midrule
Nucleon ($N$) & 939.0 & 859.0 & 91.5 \\
$\rho$-meson & 775.3 & 695.3 & 89.7 \\
$\omega$-meson & 782.6 & 702.6 & 89.8 \\
Pion ($\pi$) & 139.6 & 72.6 & 52.0 \\
Kaon ($K$) & 493.7 & 131.7 & 26.7 \\
\bottomrule
\end{tabular}
\caption{Nonperturbative mass fractions for selected hadrons. Goldstone bosons ($\pi, K$) exhibit low $f_\Omega$, while ordinary hadrons cluster at ${\sim}\,90\%$.}
\label{tab:hadrons}
\end{table}

\begin{table}[h]
\centering
\begin{tabular}{llll}
\toprule
$n_B/n_0$ & $m_\rho$ (MeV) & $m_\omega$ (MeV) & Mass drop (\%) \\
\midrule
0.0 & 775.3 & 782.6 & 0.0 \\
1.0 & 658.6 & 664.7 & $-15.0$ \\
2.0 & 595.8 & 601.2 & $-23.2$ \\
3.0 & 554.3 & 559.3 & $-28.5$ \\
\bottomrule
\end{tabular}
\caption{In-medium vector meson mass shift: a ${\sim}\,24\%$ drop of $\rho$-meson mass at $2n_0$ provides a direct signature for FAIR/HADES dielectron spectrometers.}
\label{tab:meson_shift}
\end{table}

%======================================================================
\section{Conclusion}
\label{sec:conclusion}

A phenomenological model is proposed wherein the rigidly fixed value of the nonperturbative nucleon mass component, derived from lattice QCD ($M_{\Omega,0} = 859 \pm 8$~MeV), serves as a fundamental constraint on the dense matter equation of state. The concept of an in-medium modified vacuum mass $M_\Omega(n_B)$ provides a low-parameter and strictly falsifiable method to integrate the QCD energy budget into neutron star phenomenology.

The overly stiff EOS problem is alleviated by a phase transition, and causality is ensured by a density-dependent saturating vector field. The model offers concrete observational signatures for gravitational interferometers (LIGO/Virgo) and heavy-ion physics (FAIR/HADES).

%======================================================================
\begin{thebibliography}{19}

\bibitem{Ji1995}
X.~Ji,
\emph{Breakup of hadron masses and energy-momentum tensor of QCD},
Phys.\ Rev.\ Lett.\ \textbf{74}, 1071 (1995).

\bibitem{SVZ1979}
M.\,A.~Shifman, A.\,I.~Vainshtein, V.\,I.~Zakharov,
\emph{QCD and resonance physics},
Nucl.\ Phys.\ B \textbf{147}, 385 (1979).

\bibitem{Chodos1974}
A.~Chodos et al.,
\emph{New extended model of hadrons},
Phys.\ Rev.\ D \textbf{9}, 3471 (1974).

\bibitem{Agadjanov2023}
A.~Agadjanov et al.,
\emph{Nucleon Sigma Terms},
Phys.\ Rev.\ Lett.\ \textbf{131}, 261902 (2023).

\bibitem{Gupta2021}
R.~Gupta et al.,
\emph{Pion-nucleon sigma term from lattice QCD},
Phys.\ Rev.\ Lett.\ \textbf{127}, 242002 (2021).

\bibitem{PDG2024}
S.~Navas et al.\ (PDG),
\emph{Review of Particle Physics},
Phys.\ Rev.\ D \textbf{110}, 030001 (2024).

\bibitem{Cassini2003}
B.~Bertotti, L.~Iess, P.~Tortora,
Nature \textbf{425}, 374 (2003).

\bibitem{Creminelli2017}
P.~Creminelli, F.~Vernizzi,
Phys.\ Rev.\ Lett.\ \textbf{119}, 251302 (2017).

\bibitem{LVK2024}
LIGO-Virgo-KAGRA Collaboration,
\emph{Testing GR with O4a data} (2024).

\bibitem{Will2014}
C.\,M.~Will,
\emph{The Confrontation between GR and Experiment},
Living Rev.\ Rel.\ \textbf{17}, 4 (2014).

\bibitem{Bardeen1968}
J.\,M.~Bardeen,
\emph{Non-singular general-relativistic gravitational collapse},
in Proc.\ GR5 (1968).

\bibitem{Dymnikova1992}
I.~Dymnikova,
\emph{Vacuum nonsingular black hole},
Gen.\ Rel.\ Grav.\ \textbf{24}, 235 (1992).

\bibitem{Hayward2006}
S.\,A.~Hayward,
\emph{Formation and evaporation of nonsingular black holes},
Phys.\ Rev.\ Lett.\ \textbf{96}, 031103 (2006).

\bibitem{Simpson2019}
A.~Simpson, M.~Visser,
\emph{Black-bounce to traversable wormhole},
JCAP \textbf{02}, 042 (2019).

\bibitem{Walecka1974}
J.\,D.~Walecka,
\emph{A theory of highly condensed matter},
Ann.\ Phys.\ \textbf{83}, 491 (1974).

\bibitem{Serot1986}
B.\,D.~Serot, J.\,D.~Walecka,
\emph{The relativistic nuclear many-body problem},
Adv.\ Nucl.\ Phys.\ \textbf{16}, 1 (1986).

\bibitem{BrownRho1991}
G.\,E.~Brown, M.~Rho,
\emph{Scaling effective Lagrangians in a dense medium},
Phys.\ Rev.\ Lett.\ \textbf{66}, 2720 (1991).

\bibitem{Lattimer2016}
J.\,M.~Lattimer, M.~Prakash,
\emph{The equation of state of hot, dense matter and neutron stars},
Phys.\ Rep.\ \textbf{621}, 127 (2016).

\bibitem{Drischler2021}
C.~Drischler et al.,
\emph{Chiral EFT and nuclear forces},
Prog.\ Part.\ Nucl.\ Phys.\ \textbf{121}, 103888 (2021).

\end{thebibliography}

\end{document}
